\documentclass[10pt]{article}
\usepackage[utf8]{inputenc}
\usepackage[T1]{fontenc}
\usepackage{amsmath}
\usepackage{amsfonts}
\usepackage{amssymb}
\usepackage[version=4]{mhchem}
\usepackage{stmaryrd}
\usepackage{bbold}

\begin{document}
$L / 2$\\
1

\section*{L12: The Legendre Symbol}
Definition: Let $p$ be an odd prime number and let $a \in \mathbb{Z}$ with $p+a$. The legendre symbol, denoted $\left(\frac{a}{p}\right)$, is\\
$\left(\frac{q}{p}\right):=\{1$ if $a$ is a quadratic vesidue mod $p \{-1$ if $a$ is a quadratic non residue. $\bmod p$\\
Example: $1,3,4,5,9$ are quadratic vesidues modulo II. Thus

$$
\left(\frac{1}{11}\right)=\left(\frac{3}{11}\right)=\left(\frac{4}{11}\right)=\left(\frac{5}{11}\right)=\left(\frac{9}{11}\right)=1 .
$$

$\cdot 2,6,7,8,10$ are quadratic non-vesidues modulo II. Thus $\left(\frac{2}{11}\right)=\left(\frac{6}{11}\right)=\left(\frac{7}{11}\right)=\left(\frac{8}{11}\right)=\left(\frac{10}{11}\right)=-1$.

\begin{itemize}
  \item The Legendre symbol ( $a / p$ ) indicates whether a is a quadratic residue modulo $p$ or not (in other words, whether $\left.x^{2} \equiv a \bmod p\right)$ is Solvable or not).
  \item The Legendre symbol does not tell you\\
the actual solutions to the congrnence, these need to be found by other means.\\
Example: Find the value of $\left(\frac{3}{7}\right)$. The Legendre symbol $(3 / 7)$ asks the question: is the congruence $x^{2} \equiv 3(\bmod 7)$ Solvable? We have\\
$1^{2} \equiv 1(\bmod 7) \Rightarrow 1$ is a quadratic vesidue $(\bmod 7)$\\
$2^{2} \equiv 4(\bmod 7) \Rightarrow 4$ is a quadratic vesidue $(\bmod 7)$\\
$3^{2} \equiv 2(\bmod 7) \Rightarrow 2$ is a quadratic vesidue $(\bmod 7)$\\
By Proposition 4.3 , there are exactly\\
$\frac{p-1}{2}=\frac{7-1}{2}=3$ incongrment quadratic residues\\
modulo 7. Since 1,2,4 are quadratic residues modulo 7 , then $3,5,6$ must be quadratic non-residues modulo 7 . Thus $x^{2} \equiv 3(\bmod 7)$ is not solvable, and so $\left(\frac{3}{7}\right)=-1$.
  \item The only method covered so far to compute $\left(\frac{a}{p}\right)$ is the computation of $\frac{p-1}{2}$ Congruences. We will improve our tool kit!
\end{itemize}

Theorem 4.4(Euler's Criterion): Let $p$ be an odd prime number and let $a \in \mathbb{Z}$ with $p \times a$. Then

$$
\left(\frac{a}{p}\right) \equiv a^{(p-1) / 2}(\bmod p) .
$$

Proof: Assume that $\left(\frac{a}{p}\right)=1$. Then $x^{2} \equiv a$ (nodp) has a solution, say $x_{0}$. We have

$$
\begin{aligned}
a^{(p-1) / 2} & \equiv\left(x_{0}^{2}\right)^{(p-1) / 2}(\bmod p) \\
& \equiv x_{0}^{p-1}(\bmod p) \\
& \equiv 1(\bmod p) \text { by Fermat's Little } \\
& \quad \text { Theorem } \\
& \equiv\left(\frac{q}{p}\right)(\bmod p) . \\
& \equiv 1(\bmod p) .
\end{aligned}
$$

Thus the desired result holds if $\left(\frac{q}{p}\right)=I$.\\
Now assume that $\left(\frac{a}{p}\right)=-1$. For each $i \in \mathbb{Z}$ with $1 \leq i \leq p-1$, there is a unique $j \in \mathbb{Z}$ with $1 \leqslant j \leqslant p-1$ such that $i j \equiv a(\bmod p)$\\
by Theorem 2.6. Furthermore, $i \neq j$, since\\
if $i=j$, then $i^{2} \equiv a(\bmod p)$ implies that $a$ is a quadratic residue modulo $p$ and we obtain $\left(\frac{a}{p}\right)=1$, a contradiction! Thus we can pair the integers into $1,2, \ldots, p-1$ into $\frac{p-1}{2}$ pairs, each with product a modulo $p$. Then

$$
(p-1)!\equiv a^{(p-1) / 2}(\bmod p) .
$$

Thus Wilson's Theorem implies that

$$
a^{(p-1) / 2} \equiv-1(\bmod p)
$$

from which we obtain

$$
\left(\frac{a}{p}\right) \equiv a^{(p-1) / 2}(\bmod p),
$$

and the desired result holds if $\left(\frac{a}{p}\right)=-1$.\\
Example: Find $\left(\frac{3}{7}\right)$ using Euler's criterion.\\
We have $\left(\frac{3}{7}\right)=3^{(7-1) / 2}(\bmod 7)$

$$
\begin{aligned}
& \equiv 27(\bmod 7) \\
& \equiv-1(\bmod 7) .
\end{aligned}
$$

Thus $\left(\frac{3}{7}\right)=-1$.

Proposition 4.5: Let $p$ be an odd prime, and let $a, b \in \mathbb{Z}$ with $p \nmid a$ and $p+b$. Then\\
(a) $\left(\frac{a^{2}}{p}\right)=1$\\
(b) If $a \equiv b(\bmod p)$, then $\left(\frac{a}{p}\right)=\left(\frac{b}{p}\right)$\\
(c) $\left(\frac{a b}{p}\right)=\left(\frac{a}{p}\right)\left(\frac{b}{p}\right)$.

Proof: (a) The congrnence $x^{2} \equiv a^{2}(\bmod p)$ has a Solution, namely a.\\
(b) Since $a \equiv b(\bmod p)$, the quadratic congrnence to be solved $x^{2} \equiv a \equiv b(\bmod p)$. Clearly\\
$x^{2} \equiv a(\bmod p)$ is solvable if and only if $x^{2} \equiv b(\bmod p)$ is solvable.\\
(c) We have $(p-1) / 2$

$$
\begin{aligned}
\left(\frac{a b}{p}\right) & \equiv(a b)(\bmod p) \text { by Euler's criterion } \\
& \equiv a^{(p-1) / 2} b^{(p-1) / 2}(\bmod p) \\
& \equiv\left(\frac{a}{p}\right)\left(\frac{b}{p}\right)(\bmod p) .
\end{aligned}
$$

Since the only values of the Legendre Symbol\\
are $\pm 1$, we have

$$
\left(\frac{a b}{p}\right)=\left(\frac{a}{p}\right)\left(\frac{b}{p}\right) .
$$

Example: Find $\left(\frac{-11}{7}\right)$. We have\\
$\left(\frac{-11}{7}\right)=\left(\frac{-1}{7}\right)\left(\frac{11}{7}\right)=\left(\frac{-1}{7}\right)\left(\frac{4}{7}\right)=\left(\frac{-1}{7}\right)$\\
$=\left(\frac{6}{7}\right)=-1$, where the final equality\\
follows from the previous example.

\begin{itemize}
  \item Further properties of Legendre symbols: let\\
$p$ be an odd prime and $a \in \mathbb{Z}$ with $p \times a$.\\
Then $a= \pm 2^{a_{0}} p_{1}^{a_{1}} p_{2}^{a_{2}} \cdots p_{r}^{a_{r}}$ with $p_{1}, p_{2}, \cdots, p_{r}$ odd primes distinct from $p$ and $a_{1}, a_{2}, \ldots, a_{r}$ non-negative integers. By Roposition 4.5 (c) we have that
\end{itemize}

$$
\left(\frac{q}{p}\right)=\left(\frac{ \pm 1}{p}\right)\left(\frac{2}{p}\right)^{a_{0}}\left(\frac{p 1}{p}\right)^{a_{1}}\left(\frac{p_{2}}{p}\right)^{a_{2}} \cdots\left(\frac{p r}{p}\right)^{a_{r}}
$$

Thus it suffices to compute $\left(\frac{ \pm 1}{p}\right),\left(\frac{2}{p}\right)$, and $\left(\frac{q}{p}\right)$\\
for $q$ an odd prime distinct from $p$.

Plan: $\cdot\left(\frac{1}{P}\right)=1$ by Proposition 4.5 (a).

\begin{itemize}
  \item Next will compute $\left(\frac{-1}{p}\right)$
  \item Consider $\left(\frac{2}{p}\right)$ and $\left(\frac{q}{p}\right)$ next lecture.
\end{itemize}

Theorem 4.6: Let $p$ be an odd prime. Then

$$
\left(\frac{-1}{p}\right)=(-1)^{(p-1) / 2}=\left\{\begin{array}{cl}
1 & \text { if } p \equiv 1(\bmod 4) \\
-1 & \text { if } p \equiv 3(\bmod 4)
\end{array}\right.
$$

Proof: The first equality follows from Euler's criterion with $a=-1$. To prove the second equality set $p=4 k+1$ (resp $p=4 k+3$ ) for Some $k \in \mathbb{Z}$, and compute $(-1)^{(p-1) / 2}$.

Example: $\left(\frac{-1}{7}\right)=-1$ by Theorem 4.6 Since $7 \equiv 3(\bmod 4)$.

\begin{itemize}
  \item $\left(\frac{-1}{53}\right)=1$ by Theorem 4.6 since $53 \equiv 1$ (nod) 4 .
\end{itemize}

Lemma 4.7 (Gauss Lemma): Let $p$ be an odd prime and let $a \in \mathbb{Z}$ with $p \nless a$. Let $n$ be the number of least positive residnes of the integers $a, 2 a, 3 a, \ldots,\left(\frac{p-1}{2}\right) a$ that are greater than $P / 2$. Then

$$
\left(\frac{a}{p}\right)=(-1)^{n} .
$$

Example: Use Ganss Lemma to find $\left(\frac{6}{11}\right)$. We have $\left(\frac{6}{11}\right)=(-1)^{n}$ where $n$ is the number of least positive residues of the integers:

$$
6,2 \times 6,3 \times 6,4 \times 6,5 \times 6
$$

that are greater than $11 / 2=5.5$. We have $6 \equiv 6(\bmod 11)$

Thus\\
$2 \times 6 \equiv 12 \equiv 1(\bmod 11)$\\
$3 \times 6 \equiv 18 \equiv 7(\bmod 11)$\\
$\left(\frac{6}{11}\right)=(-1)^{3}=-1$\\
$4 \times 6 \equiv 24 \equiv 2$ (nod 11) by Gauss\\
$5 \times 6 \equiv 30 \equiv 8$ (mod 11) Lemma.


\end{document}